\section{Concert Tickets}
\label{sec:g-concert}
\source{CSES 1091 (Notion: Greedy)}{Easy--Medium}

\problemstatement{There are $n$ tickets with prices $h_i$. Customers arrive
one by one; customer $j$ buys the most expensive remaining ticket whose price
does not exceed $t_j$, or gets nothing ($-1$). Print what each customer pays
($n,m\le2\cdot10^5$).}

\begin{patternbox}
``Largest remaining element $\le x$, then delete it'' $\Rightarrow$
\texttt{multiset} with \texttt{upper\_bound} and one step back.
\end{patternbox}

\subsection*{Approach and intuition}

The rule is given by the statement, so the whole task is choosing a data
structure supporting \emph{predecessor query} and \emph{delete} in
$O(\log n)$ with duplicates: a \texttt{std::multiset}.
\texttt{upper\_bound(t)} returns the first price $>t$; the element just
before it (if any) is the largest price $\le t$.

\subsection*{Algorithm}
\begin{algorithmic}[1]
\State insert all prices into multiset $S$
\For{each budget $t$}
  \State $it\gets S.\text{upper\_bound}(t)$
  \If{$it=S.\text{begin}()$} \State print $-1$
  \Else \State $it\gets it-1$; print $*it$; erase $it$ \EndIf
\EndFor
\end{algorithmic}

\dryrun{Prices $\{3,5,5,7,8\}$. $t=4$: pays $3$. $t=8$: pays $8$. $t=3$:
nothing $\le3$ left, $-1$. \checkmark}

\subsection*{C++17 implementation}
\cppfile{greedy/12_concert_tickets.cpp}

\begin{complexitybox}
\textbf{Time:} $O((n+m)\log n)$. \textbf{Space:} $O(n)$.
\end{complexitybox}

\begin{pitfallbox}
\begin{itemize}
  \item \texttt{S.erase(value)} would delete \emph{all} copies of that
        price; erase the iterator.
  \item \texttt{lower\_bound(t)} is off by one for equal prices.
\end{itemize}
\end{pitfallbox}
